% SPDX-License-Identifier: CC-BY-4.0
% Copyright (C) 2026 Fabian Januszewski, Abhijit (the Langlandscape team)
% GENERATED by tools/extract_statements.py from mathematics/chapters/ -- do not edit by hand.

\begin{quoted}{Theorem}{7.7}{54}{thm:s02-SA}{strong approximation for $B^1$ at $S=\{\infty,q\}$}{chapters/s02/2-quotient.tex}
Since $q\nmid D$, $B_q\cong M_2(\Qq)$ is split and $B^1_q\cong\SL_2(\Qq)$ is non-compact. Hence $B^1$ (embedded diagonally) is
dense in the restricted product $\Bhat^{1,(q)}:=\prod'_{\ell\ne q}B^1_\ell$ (with respect to the $\cO^1_\ell$).
\smallskip

\noindent\emph{Consequence (the shape used; \texttt{if:s02-SA}).} For every $q$-admissible $U$,
\[\Bhat^1=B^1\big(B^1_q\times U^{1,(q)}\big).\tag{SA$_U$}\]
\par\smallskip\noindent{\scriptsize\itshape [Status paragraph following the statement in the chapter:]}\par
\noindent\emph{Status.} The theorem is \stcited{}:
\lit{voight\_2021\_quaternion-algebras}{Main Thm.~28.5.3, p.~485}, applied to the global ring $R=\Z[1/q]$ with eligible set
$S=\{\infty,q\}$. In Voight's notation $\Bhat=\Bhat^S=\prod'_{v\notin S}B_v$ is the $S$-finite adele ring
(\lit{voight\_2021\_quaternion-algebras}{27.6.4, p.~469}; there, when $S$ is clear from the context, the superscript $S$ is dropped and replaced by hats), so his
$\Bhat^1$ is our $\Bhat^{1,(q)}$. The hypothesis ``$B$ is $S$-indefinite'' (\lit{voight\_2021\_quaternion-algebras}{Def.~28.5.1,
p.~485}: $S$ contains a place unramified in $B$) holds because $q\in S$ and $q\nmid D$.
Equivalently \lit{vigneras\_1980\_arithmetique-des-algebres-de-quaternions}{Ch.~III, Thm.~4.3, p.~81} ($H^1_KH^1_S$ dense in
$H^1_{\A}$ if $H^1_S$ is not compact; page-checked) and
\lit{platonov-rapinchuk\_1994\_algebraic-groups-and-number-theory\_2up}{\S7.4, Thm.~7.12, p.~427} (page-checked in the M8 litnote).
Kneser \lit{kneser\_1969\_lectures-on-galois-cohomology-of-classical-groups}{\S5.5, Sublemma~b, p.~94 (Tata pagination; p.~72 of the scan)} only states the result;
it is listed for orientation and not relied upon. The consequence (SA$_U$) is \stproved{} below from the theorem. It is
\lit{voight\_2021\_quaternion-algebras}{28.4.8, p.~484} (density implies $\Bhat^1\subseteq B^1\hat\cO^1$) and
\lit{voight\_2021\_quaternion-algebras}{Cor.~28.5.12, p.~486}; the shape (SA$_U$), with the factor $B^1_q$, is the projection away from $\infty$ of the \emph{second}
equality of (28.5.13) there (which multiplies in $B^1_S=B^1_\infty\times B^1_q$). Both are stated only for the norm-one units $\hat\cO^1$ of an
$R$-order $\cO$, not for a general open subgroup; the same proof extends, since it uses only that $\hat\cO^1$ is open (proof below); this matters because $U^{(q)}$ need not be the unit group of an order (e.g.\ $\Uone{N}$-levels).
In Lean, (SA$_U$) for all admissible $U$ is registered as the hypothesis $\HypSA$, \lean{Langlandscape.Hyp.SA} (shape
\lean{Langlandscape.Tree.SAat}, quantified over \pathtt{Tree.IsQAdmissible}); this is a registered hypothesis, not a proof, and the
theorem stays \stcited. The derivation of (SA$_U$) from density below is not formalised. Per level, $\HypSA$ is replaced by
Prop.~7.9.
\end{quoted}

\begin{quoted}{Proposition}{7.9}{55}{prop:s02-BG}{SA for one level as a finite check}{chapters/s02/2-quotient.tex}
Let $U$ be $q$-admissible. The maps
\begin{gather*}
\phi_0:\GU\backslash V_0(\cT_q)\to\XB U,\ \GU gv_0\mapsto[g];\\
\phi_1:\GU\backslash V_1(\cT_q)\to\XB U,\ \GU g\alpha v_0\mapsto[g\alpha];\\
\phi_E:\GU\backslash E(\cT_q)\to\XB{U'},\ \GU ge_0\mapsto[g]
\end{gather*}
($g\in B^1_q$) are well defined and \emph{injective}, and $\phi=(\phi_0,\phi_1,\phi_E):\GU\backslash\cT_q\to\BG(U,q)$ is a locally
surjective morphism of graphs. The following are equivalent:
\begin{enumerate}[label=(\roman*)]
\item $\phi_0$ is surjective;
\item (SA$_U$) holds and $\Xone(U)\to\XB{U}$ is surjective;
\item $\BG(U,q)$ is connected.
\end{enumerate}
Moreover:
\begin{enumerate}[label=(\alph*)]
\item (i) implies that $\phi_0,\phi_1,\phi_E$ are all bijective, so $\phi$ is an isomorphism of graphs $\GU\backslash\cT_q\cong\BG(U,q)$.
  This implication does not use the connectedness of $\cT_q$.
\item In all cases, $\phi$ is an isomorphism of $\GU\backslash\cT_q$ onto a union of connected components of $\BG(U,q)$ containing
  $([1],0)$; using that $\cT_q$ is connected, onto the connected component of $([1],0)$.
\end{enumerate}
\par\smallskip\noindent{\scriptsize\itshape [Status paragraph following the statement in the chapter:]}\par
\noindent\emph{Status.} \stproved{} (new in this plan). No strong approximation is used. The connectedness of $\cT_q$
(Lemma~7.3(a), \stcited) enters only in (i)$\Rightarrow$(iii) and in the last clause of (b); the implications
(iii)$\Rightarrow$(i)$\Leftrightarrow$(ii), the injectivity, local surjectivity, and (a) are tree-free.
In particular Thm.~7.7 together with Hasse--Schilling (Lemma~7.6(b)) implies (ii), hence (i) and the
bijections of (a) (tree-free), and hence (iii) (using connectedness of $\cT_q$). As hypotheses, SA alone does not give (ii).
Conversely, for an explicit $(B,U,q)$ the condition (iii) is a finite computation, and it discharges (SA$_U$) and the surjectivity
of $\Xone(U)\to\XB U$ for that level. Classically, SA enters at this point through the definite analogue of Ihara's lemma:
\lit{diamond-taylor\_1994\_non-optimal-levels-of-mod-l-modular-representations}{Lemma~2 and Cor.~1, pp.~445--446} show that $T_\ell-2$, $S_\ell-1$ annihilate $A[l]$, $A$ the kernel of
$(f_1,f_2)\mapsto f_1+\varphi_pf_2$ on $S_k(U;\Q_l/\Z_l)^2$, for almost all $\ell\equiv1\bmod d$ (Lemma~2; hypotheses $p\nmid N_U l$), whence the Eisenstein statement (Cor.~1), and the paper says this lemma ``plays the role of Ihara's result'' (p.~445).
\emph{Lean status.} \stformal{} (all tree-free parts; in the coset model of Lemma~7.3, with $\phi_0,\phi_1$ combined
into one map \texttt{phiV} on all vertex orbits to $\XB U\sqcup\XB U$):
\lean{Langlandscape.Tree.phiV\_injective}, \lean{Langlandscape.Tree.phiE\_injective},
\lean{Langlandscape.Tree.phiV\_oO}, \lean{Langlandscape.Tree.phiV\_tO} (Step~3),
\lean{Langlandscape.Tree.phi\_locSurj\_type0}, \lean{Langlandscape.Tree.phi\_locSurj\_type1} (Step~4);
(i) is \lean{Langlandscape.Tree.surjType0} (every class of $\XB U$ is $[g]$ with $g\in\SL_2(\Qq)$), and
\lean{Langlandscape.Tree.surjType0\_iff} identifies it with the surjectivity of $\phi_0$;
\lean{Langlandscape.Tree.surjType0\_of\_BGConnected} ((iii)$\Rightarrow$(i)),
\lean{Langlandscape.Tree.bijective\_of\_BGConnected} ((iii)$\Rightarrow$ bijections),
\lean{Langlandscape.Tree.bijective\_of\_surjType0} ((a)), \lean{Langlandscape.Tree.surj\_of\_SA} ((ii)$\Rightarrow$(i)),
\lean{Langlandscape.Tree.SA\_of\_surj} ((i)$\Rightarrow$(ii)). None of these takes a \texttt{Hyp.*} argument.
(i)$\Rightarrow$(iii) is \lean{Langlandscape.Tree.BGConnected\_of\_surj}, formalized modulo \texttt{Hyp.BTtree}; the ``in particular''
is \lean{Langlandscape.Tree.BGConnected\_of\_SA\_HS}, formalized modulo $\HypSA$, $\HypHS$, \texttt{Hyp.BTtree} (and
\lean{Langlandscape.Tree.surjType0\_of\_SA\_HS} for (i), modulo $\HypSA$, $\HypHS$ only). The last clause of (b) (image $=$ the
component of $([1],0)$) also uses connectedness of $\cT_q$; it has no separate Lean declaration (not formalised; it would take
\texttt{Hyp.BTtree}).
\emph{Hypothesis on $U$ in Lean} (audit note): all of these except \pathtt{BGConnected\_of\_SA\_HS} and \pathtt{surjType0\_of\_SA\_HS}
assume only that $U$ is a $q$-level (\texttt{IsQLevel}: $U\le\hat\cO^\times$ of finite index, factoring at $q$), not that $U$ is
$q$-admissible; the condition $\Nrd U=\Zhat^\times$ is not needed for the proposition, in line with the proof above (Step~2 uses only
$\Nrd U\subseteq\Zhat^\times$).
\emph{Deviation from the plan tex (refinement).} The tex's ``Moreover, in all cases the maps are injective and $\phi$ is an isomorphism
onto the component of $(1,0)$'' is split here into the injectivity/local surjectivity in the statement, (a) ((i)$\Rightarrow$ bijections,
tree-free) and (b) (union of components, tree-free; single component, with the tree).
\end{quoted}

\begin{quoted}{Proposition}{7.10}{57}{prop:s02-quot}{quotient graph; seed Prop.~quot}{chapters/s02/2-quotient.tex}
Let $U$ be $q$-admissible and assume (iii) of Prop.~7.9. By Prop.~7.9 this holds given
Thm.~7.7, Hasse--Schilling \emph{and} the connectedness of $\cT_q$ (Lemma~7.3(a), \stcited; \textbf{[tree]}):
$\HypSA\wedge\HypHS\wedge\mathrm{Hyp.BTtree}\Rightarrow$ (iii) (Lean \pathtt{Tree.BGConnected\_of\_SA\_HS}, which takes \texttt{Hyp.BTtree});
the bijections of (a) need only $\HypSA\wedge\HypHS$. As hypotheses, SA alone does not suffice.
(iii) itself is the working hypothesis here.
\begin{enumerate}[label=(\alph*)]
\item There are bijections
  \begin{gather*}
    \GU\backslash\{\text{type-}0\text{ vertices}\}\ \xrightarrow{\ gv_0\mapsto[g]\ }\ \XB{U},\\
    \GU\backslash\{\text{type-}1\text{ vertices}\}\ \xrightarrow{\ g\alpha v_0\mapsto[g\alpha]\ }\ \XB{U},\\
    \GU\backslash\{\text{edges}\}\ \xrightarrow{\ ge_0\mapsto[g]\ }\ \XB{\Uzero{q}U}
  \end{gather*}
  ($g\in B^1_q$). Under these bijections the endpoint maps $o,t$ become $s_0,s_1$. In other words,
  $\GU\backslash\cT_q\cong\BG(U,q)$ as graphs.
\item For $g\in B^1_q$, $\Stab_{\GU}(gv_0)=\Gamma_g$ (level $U$), $\Stab_{\GU}(g\alpha v_0)=\Gamma_{g\alpha}$
  (level $U$) and $\Stab_{\GU}(ge_0)=\Gamma_g$ (level $\Uzero{q}U$). For another representative
  $x'=\gamma xu$ one has $\Gamma_{x'}=\gamma\Gamma_x\gamma^{-1}$. Each $\Gamma_x$ has norm $1$
  (Lemma~1.7(a)), so it lies in $\GU$ after conjugation.
\item Let $T$ be any Hecke operator supported away from $q$ ($T_\ell,S_\ell,U_\ell$ for $\ell\mid M$,
  $\langle d\rangle$, $U_p$ with the normalisation $\NUp$ of stage 01, and $\Pi_\ell$ for $\ell\mid D$),
  acting compatibly on levels $U$ and $\Uzero{q}U$. Then $s_0^*,s_1^*,s_{0*},s_{1*}$ commute with $T$ on
  Brandt modules $\Brandt{\cdot}{V}$ and on their homological versions (on $\Brandt{\cdot}{V}$, $s_{j*}$ means the trace
  $\mathrm{tr}_j$ of Def.~7.13).
\end{enumerate}
\par\smallskip\noindent{\scriptsize\itshape [Status paragraph following the statement in the chapter:]}\par
\noindent\emph{Status.} (a): \stproved{} given (iii); \stpmc{} via Thm.~7.7 and Hasse--Schilling (through
(ii)$\Rightarrow$(i)$\Rightarrow$(a), which does not use connectedness of the tree). (b): \stproved; \emph{refinement of the plan tex} (which has ``given (iii)''): the three stabiliser identities
hold for every $g\in B^1_q$ without (iii); (iii) is used only to know that every class of $\XB U$, $\XB{U'}$ has such a tree
representative. (c): \stproved{} (unconditional).
\end{quoted}

\begin{quoted}{Proposition}{7.19}{66}{prop:s02-harm}{harmonic cochains; seed Prop.~harm}{chapters/s02/3-cohomology.tex}
Under the hypotheses of Prop.~7.18(b), let $c$ be a function on oriented edges with values in $V$
such that $c(\gamma e)=\gamma_p\,c(e)$ for $\gamma\in\GU$, $c(\bar e)=-c(e)$, and $\sum_{o(e)=v}c(e)=0$ for all
$v$. The $\GU$-invariant $V$-valued harmonic cochains of this kind are
\[\ker\Big(\Brandt{\Uzero{q}U}{V}\xrightarrow{(\mathrm{tr}_0,\,\mathrm{tr}_1)}\Brandt{U}{V}^{2}\Big)\ \cong\ H_1(\GU,V),\]
with the trace maps $\mathrm{tr}_j$ (the plan's ``$s_{j*}$'') of Def.~7.13. If $-1\in U$, assume $-1$ acts trivially on $V$
(the standing hypothesis of Def.~7.13, needed for $\nu$); otherwise replace $V$ by $V^+$, as $V^-$ contributes $0$ to both sides
(Remark~7.15). This is the \emph{$q$-new} submodule of the Brandt
module of level $\Uzero{q}U$; the term is defined here as this kernel. The isomorphism is $\nu_{U'}^{-1}$ followed by
Prop.~7.18(b), $H_1(\GU,V)=\ker\partial$.
\par\smallskip\noindent{\scriptsize\itshape [Status paragraph following the statement in the chapter:]}\par
\noindent The description of the harmonic cochains as $\ker(\mathrm{tr}_0,\mathrm{tr}_1)$ is \stproved{} given (iii); the isomorphism with
$H_1(\GU,V)$ is \stpmc{} given (iii) (via Prop.~7.18(b); registered in Lean as hypothesis \pathtt{Hyp.SerreCoh}). (In
the sum, ``$o(e)=v$'' runs over all oriented edges with origin $v$, in both orientation classes.)
\emph{Lean status.} For trivial coefficients $V=R$ the $q$-new module $\ker(\mathrm{tr}_0,\mathrm{tr}_1)$ is
\lean{Langlandscape.Tree.qNew}, and $\nu_{U'}:\ker\partial\cong\ker(\mathrm{tr}_0,\mathrm{tr}_1)$ is \stformal{}
\lean{Langlandscape.Tree.qNewEquivKerPartial} (also \lean{Langlandscape.Tree.normMap\_restrict\_bijective}), assuming all $w_x$
at level $U$ are units in $R$ (remark below); no (iii) is needed for it. The identification with harmonic cochains on the tree is not
formalised.
\end{quoted}
